I can translate the Japanese comments in the given TeX code:

\begin{lstlisting}[language=c++,label={code:25_01},caption={Approximation and Verification}]
X <- rnorm(100000,mean=0,sd=1)
length(X[X>1])/100000
# Verification
1-pnorm(1,mean=0,sd=1)
\end{lstlisting} 

The first comment reads "Approximation", which refers to the code that follows, where `rnorm` function is used to generate a random set of 100000 numbers following a normal distribution with mean 0 and standard deviation 1, stored in the variable `X`. The following line of code counts how many of these numbers are greater than 1, and then divides that value by the number of total values (100000), to get an approximation of the probability of generating a value greater than 1 from this distribution.

The second comment reads "Verification", which refers to the next line of code that calculates the true probability of generating a value greater than 1 from this distribution, using the `pnorm` function. The overall code shows an example of how to use sampling and verification to estimate probabilities of events from a given distribution.
